\section{Amenazas a la Validez}
\label{sec:amenazas-validez}

\subsection{Validez de constructo}
El p95 de latencia de k6 se usa como proxy de ``rendimiento percibido'', pero no se midió tiempo de
renderizado en el navegador ni Web Vitals bajo carga real de usuario (solo Lighthouse en condición de
reposo, sin carga concurrente). El puntaje SUS, cuando se aplique, medirá usabilidad percibida
declarada, no eficiencia real de tarea (tiempo para completar un flujo).

\subsection{Validez interna}
Las mediciones de rendimiento y Lighthouse se ejecutaron en la máquina de desarrollo personal de un
integrante del equipo, con otras aplicaciones activas (Discord, Steam, Brave), confirmado con
\texttt{Get-Process | Sort CPU}. Esto es una amenaza real y documentada, no hipotética: el First
Contentful Paint medido (2.9--5.9\,s) es alto para un bundle de $\sim$100\,KB transferido servido en
localhost sin latencia de red real, lo que sugiere contención de CPU no relacionada con la aplicación
misma, no un problema de rendimiento intrínseco del sistema.

\subsection{Validez externa}
Los 50 usuarios concurrentes de k6 son un valor fijado por el RNF-01, no derivado de datos reales de uso
institucional de la UTEQ (que no existen para este sistema, al ser nuevo). No hay garantía de que ese
sea el nivel de concurrencia real durante un periodo de matrícula/defensas masivas. El usuario de
demostración y las pruebas manuales no reemplazan una prueba piloto con usuarios reales del dominio
(estudiantes, coordinadores) --- que es precisamente la brecha que el SUS pendiente busca cerrar.

\subsection{Validez de conclusión}
El tamaño de muestra de Lighthouse ($n=3$ por perfil) es demasiado pequeño para una prueba de hipótesis
formal entre desktop y mobile (Sección~\ref{sec:evaluacion}); solo se reportan diferencias descriptivas.
La revisión de trabajos relacionados (Sección~\ref{sec:trabajos-relacionados}) no siguió un protocolo de
búsqueda sistemática formal con doble revisor, lo que limita la conclusión de que la cobertura de
literatura es exhaustiva --- se declaró explícitamente como una búsqueda dirigida, no una SLR completa.
Finalmente, el propio historial de datos fabricados detectados en este proyecto (Sección~\ref{sec:discusion},
RQ4) es una amenaza de conclusión de segundo orden: obliga a que cada cifra en este informe sea trazable
a \texttt{docs/mediciones/DATA-PROVENANCE.md}, precisamente porque la confianza por defecto en la
documentación de este proyecto específico ya se demostró insuficiente en al menos cinco episodios
distintos.
\newpage
